\IfFileExists{stacks-project.cls}{%
\documentclass{stacks-project}
}{%
\documentclass{amsart}
}

% For dealing with references we use the comment environment
\usepackage{verbatim}
\newenvironment{reference}{\comment}{\endcomment}
%\newenvironment{reference}{}{}
\newenvironment{slogan}{\comment}{\endcomment}
\newenvironment{history}{\comment}{\endcomment}

% For commutative diagrams we use Xy-pic
\usepackage[all]{xy}

% We use 2cell for 2-commutative diagrams.
\xyoption{2cell}
\UseAllTwocells

% We use multicol for the list of chapters between chapters
\usepackage{multicol}

% This is generally recommended for better output
\usepackage{lmodern}
\usepackage[T1]{fontenc}

% For cross-file-references
\usepackage{xr-hyper}

% Package for hypertext links:
\usepackage{hyperref}

% For any local file, say "hello.tex" you want to link to please
% use \externaldocument[hello-]{hello}
\externaldocument[introduction-]{introduction}
\externaldocument[conventions-]{conventions}
\externaldocument[sets-]{sets}
\externaldocument[categories-]{categories}
\externaldocument[topology-]{topology}
\externaldocument[sheaves-]{sheaves}
\externaldocument[sites-]{sites}
\externaldocument[stacks-]{stacks}
\externaldocument[fields-]{fields}
\externaldocument[algebra-]{algebra}
\externaldocument[brauer-]{brauer}
\externaldocument[homology-]{homology}
\externaldocument[derived-]{derived}
\externaldocument[simplicial-]{simplicial}
\externaldocument[more-algebra-]{more-algebra}
\externaldocument[smoothing-]{smoothing}
\externaldocument[modules-]{modules}
\externaldocument[sites-modules-]{sites-modules}
\externaldocument[injectives-]{injectives}
\externaldocument[cohomology-]{cohomology}
\externaldocument[sites-cohomology-]{sites-cohomology}
\externaldocument[dga-]{dga}
\externaldocument[dpa-]{dpa}
\externaldocument[sdga-]{sdga}
\externaldocument[hypercovering-]{hypercovering}
\externaldocument[schemes-]{schemes}
\externaldocument[constructions-]{constructions}
\externaldocument[properties-]{properties}
\externaldocument[morphisms-]{morphisms}
\externaldocument[coherent-]{coherent}
\externaldocument[divisors-]{divisors}
\externaldocument[limits-]{limits}
\externaldocument[varieties-]{varieties}
\externaldocument[topologies-]{topologies}
\externaldocument[descent-]{descent}
\externaldocument[perfect-]{perfect}
\externaldocument[more-morphisms-]{more-morphisms}
\externaldocument[flat-]{flat}
\externaldocument[groupoids-]{groupoids}
\externaldocument[more-groupoids-]{more-groupoids}
\externaldocument[etale-]{etale}
\externaldocument[chow-]{chow}
\externaldocument[intersection-]{intersection}
\externaldocument[pic-]{pic}
\externaldocument[weil-]{weil}
\externaldocument[adequate-]{adequate}
\externaldocument[dualizing-]{dualizing}
\externaldocument[duality-]{duality}
\externaldocument[discriminant-]{discriminant}
\externaldocument[derham-]{derham}
\externaldocument[local-cohomology-]{local-cohomology}
\externaldocument[algebraization-]{algebraization}
\externaldocument[curves-]{curves}
\externaldocument[resolve-]{resolve}
\externaldocument[models-]{models}
\externaldocument[functors-]{functors}
\externaldocument[equiv-]{equiv}
\externaldocument[pione-]{pione}
\externaldocument[etale-cohomology-]{etale-cohomology}
\externaldocument[proetale-]{proetale}
\externaldocument[relative-cycles-]{relative-cycles}
\externaldocument[more-etale-]{more-etale}
\externaldocument[trace-]{trace}
\externaldocument[crystalline-]{crystalline}
\externaldocument[spaces-]{spaces}
\externaldocument[spaces-properties-]{spaces-properties}
\externaldocument[spaces-morphisms-]{spaces-morphisms}
\externaldocument[decent-spaces-]{decent-spaces}
\externaldocument[spaces-cohomology-]{spaces-cohomology}
\externaldocument[spaces-limits-]{spaces-limits}
\externaldocument[spaces-divisors-]{spaces-divisors}
\externaldocument[spaces-over-fields-]{spaces-over-fields}
\externaldocument[spaces-topologies-]{spaces-topologies}
\externaldocument[spaces-descent-]{spaces-descent}
\externaldocument[spaces-perfect-]{spaces-perfect}
\externaldocument[spaces-more-morphisms-]{spaces-more-morphisms}
\externaldocument[spaces-flat-]{spaces-flat}
\externaldocument[spaces-groupoids-]{spaces-groupoids}
\externaldocument[spaces-more-groupoids-]{spaces-more-groupoids}
\externaldocument[bootstrap-]{bootstrap}
\externaldocument[spaces-pushouts-]{spaces-pushouts}
\externaldocument[spaces-chow-]{spaces-chow}
\externaldocument[groupoids-quotients-]{groupoids-quotients}
\externaldocument[spaces-more-cohomology-]{spaces-more-cohomology}
\externaldocument[spaces-simplicial-]{spaces-simplicial}
\externaldocument[spaces-duality-]{spaces-duality}
\externaldocument[formal-spaces-]{formal-spaces}
\externaldocument[restricted-]{restricted}
\externaldocument[spaces-resolve-]{spaces-resolve}
\externaldocument[formal-defos-]{formal-defos}
\externaldocument[defos-]{defos}
\externaldocument[cotangent-]{cotangent}
\externaldocument[examples-defos-]{examples-defos}
\externaldocument[algebraic-]{algebraic}
\externaldocument[examples-stacks-]{examples-stacks}
\externaldocument[stacks-sheaves-]{stacks-sheaves}
\externaldocument[criteria-]{criteria}
\externaldocument[artin-]{artin}
\externaldocument[quot-]{quot}
\externaldocument[stacks-properties-]{stacks-properties}
\externaldocument[stacks-morphisms-]{stacks-morphisms}
\externaldocument[stacks-limits-]{stacks-limits}
\externaldocument[stacks-cohomology-]{stacks-cohomology}
\externaldocument[stacks-perfect-]{stacks-perfect}
\externaldocument[stacks-introduction-]{stacks-introduction}
\externaldocument[stacks-more-morphisms-]{stacks-more-morphisms}
\externaldocument[stacks-geometry-]{stacks-geometry}
\externaldocument[moduli-]{moduli}
\externaldocument[moduli-curves-]{moduli-curves}
\externaldocument[examples-]{examples}
\externaldocument[exercises-]{exercises}
\externaldocument[guide-]{guide}
\externaldocument[desirables-]{desirables}
\externaldocument[coding-]{coding}
\externaldocument[obsolete-]{obsolete}
\externaldocument[fdl-]{fdl}
\externaldocument[index-]{index}

% Theorem environments.
%
\theoremstyle{plain}
\newtheorem{theorem}[subsection]{Theorem}
\newtheorem{proposition}[subsection]{Proposition}
\newtheorem{lemma}[subsection]{Lemma}

\theoremstyle{definition}
\newtheorem{definition}[subsection]{Definition}
\newtheorem{example}[subsection]{Example}
\newtheorem{exercise}[subsection]{Exercise}
\newtheorem{situation}[subsection]{Situation}

\theoremstyle{remark}
\newtheorem{remark}[subsection]{Remark}
\newtheorem{remarks}[subsection]{Remarks}

\numberwithin{equation}{subsection}

% Macros
%
\def\lim{\mathop{\mathrm{lim}}\nolimits}
\def\colim{\mathop{\mathrm{colim}}\nolimits}
\def\Spec{\mathop{\mathrm{Spec}}}
\def\Hom{\mathop{\mathrm{Hom}}\nolimits}
\def\Ext{\mathop{\mathrm{Ext}}\nolimits}
\def\SheafHom{\mathop{\mathcal{H}\!\mathit{om}}\nolimits}
\def\SheafExt{\mathop{\mathcal{E}\!\mathit{xt}}\nolimits}
\def\Sch{\mathit{Sch}}
\def\Mor{\mathop{\mathrm{Mor}}\nolimits}
\def\Ob{\mathop{\mathrm{Ob}}\nolimits}
\def\Sh{\mathop{\mathit{Sh}}\nolimits}
\def\NL{\mathop{N\!L}\nolimits}
\def\CH{\mathop{\mathrm{CH}}\nolimits}
\def\proetale{{pro\text{-}\acute{e}tale}}
\def\etale{{\acute{e}tale}}
\def\QCoh{\mathit{QCoh}}
\def\Ker{\mathop{\mathrm{Ker}}}
\def\Im{\mathop{\mathrm{Im}}}
\def\Coker{\mathop{\mathrm{Coker}}}
\def\Coim{\mathop{\mathrm{Coim}}}

% Boxtimes
%
\DeclareMathSymbol{\boxtimes}{\mathbin}{AMSa}{"02}

%
% Macros for moduli stacks/spaces
%
\def\QCohstack{\mathcal{QC}\!\mathit{oh}}
\def\Cohstack{\mathcal{C}\!\mathit{oh}}
\def\Spacesstack{\mathcal{S}\!\mathit{paces}}
\def\Quotfunctor{\mathrm{Quot}}
\def\Hilbfunctor{\mathrm{Hilb}}
\def\Curvesstack{\mathcal{C}\!\mathit{urves}}
\def\Polarizedstack{\mathcal{P}\!\mathit{olarized}}
\def\Complexesstack{\mathcal{C}\!\mathit{omplexes}}
% \Pic is the operator that assigns to X its picard group, usage \Pic(X)
% \Picardstack_{X/B} denotes the Picard stack of X over B
% \Picardfunctor_{X/B} denotes the Picard functor of X over B
\def\Pic{\mathop{\mathrm{Pic}}\nolimits}
\def\Picardstack{\mathcal{P}\!\mathit{ic}}
\def\Picardfunctor{\mathrm{Pic}}
\def\Deformationcategory{\mathcal{D}\!\mathit{ef}}
\usepackage{cleveref}
\crefname{originaldefinition}{original definition}{original definitions}
\Crefname{originaldefinition}{Original definition}{Original definitions}
\makeatletter
\newlabel{situation-coherent}{{08KB}{source}{}{sourcepointer-1}{}}
\newlabel{situation-coherent@cref}{{[originaldefinition][0][]08KB}{[1][0][]source}{}{}{}}
\makeatother
\crefname{originalsection}{original section}{original sections}
\Crefname{originalsection}{Original section}{Original sections}
\makeatletter
\newlabel{artin-section-tangent-spaces}{{07WY}{source}{}{sourcepointer-2}{}}
\newlabel{artin-section-tangent-spaces@cref}{{[originalsection][0][]07WY}{[1][0][]source}{}{}{}}
\makeatother
\crefname{originallemma}{original lemma}{original lemmas}
\Crefname{originallemma}{Original lemma}{Original lemmas}
\makeatletter
\newlabel{lemma-coherent-RS-star}{{08LQ}{source}{}{sourcepointer-3}{}}
\newlabel{lemma-coherent-RS-star@cref}{{[originallemma][0][]08LQ}{[1][0][]source}{}{}{}}
\makeatother
\crefname{originalsection}{original section}{original sections}
\Crefname{originalsection}{Original section}{Original sections}
\makeatletter
\newlabel{artin-section-inf}{{07Y6}{source}{}{sourcepointer-4}{}}
\newlabel{artin-section-inf@cref}{{[originalsection][0][]07Y6}{[1][0][]source}{}{}{}}
\makeatother
\crefname{originallemma}{original lemma}{original lemmas}
\Crefname{originallemma}{Original lemma}{Original lemmas}
\makeatletter
\newlabel{artin-lemma-deformation-category}{{07WU}{source}{}{sourcepointer-5}{}}
\newlabel{artin-lemma-deformation-category@cref}{{[originallemma][0][]07WU}{[1][0][]source}{}{}{}}
\makeatother
\crefname{originalremark}{original remark}{original remarks}
\Crefname{originalremark}{Original remark}{Original remarks}
\makeatletter
\newlabel{remark-coherent-base-change}{{08LP}{source}{}{sourcepointer-6}{}}
\newlabel{remark-coherent-base-change@cref}{{[originalremark][0][]08LP}{[1][0][]source}{}{}{}}
\makeatother
\crefname{originallemma}{original lemma}{original lemmas}
\Crefname{originallemma}{Original lemma}{Original lemmas}
\makeatletter
\newlabel{artin-lemma-fibre-product-tangent-spaces}{{07X2}{source}{}{sourcepointer-7}{}}
\newlabel{artin-lemma-fibre-product-tangent-spaces@cref}{{[originallemma][0][]07X2}{[1][0][]source}{}{}{}}
\makeatother
\crefname{originallemma}{original lemma}{original lemmas}
\Crefname{originallemma}{Original lemma}{Original lemmas}
\makeatletter
\newlabel{artin-lemma-finite-dimension}{{07X1}{source}{}{sourcepointer-8}{}}
\newlabel{artin-lemma-finite-dimension@cref}{{[originallemma][0][]07X1}{[1][0][]source}{}{}{}}
\makeatother
\crefname{originallemma}{original lemma}{original lemmas}
\Crefname{originallemma}{Original lemma}{Original lemmas}
\makeatletter
\newlabel{lemma-coherent-over-first-order-thickening}{{08W7}{source}{}{sourcepointer-9}{}}
\newlabel{lemma-coherent-over-first-order-thickening@cref}{{[originallemma][0][]08W7}{[1][0][]source}{}{}{}}
\makeatother
\crefname{originallemma}{original lemma}{original lemmas}
\Crefname{originallemma}{Original lemma}{Original lemmas}
\makeatletter
\newlabel{defos-lemma-flat-ringed-topoi}{{08VW}{source}{}{sourcepointer-10}{}}
\newlabel{defos-lemma-flat-ringed-topoi@cref}{{[originallemma][0][]08VW}{[1][0][]source}{}{}{}}
\makeatother
\crefname{originallemma}{original lemma}{original lemmas}
\Crefname{originallemma}{Original lemma}{Original lemmas}
\makeatletter
\newlabel{spaces-perfect-lemma-ext-finite}{{0D0T}{source}{}{sourcepointer-11}{}}
\newlabel{spaces-perfect-lemma-ext-finite@cref}{{[originallemma][0][]0D0T}{[1][0][]source}{}{}{}}
\makeatother
\title{Finite dimensionality of tangent and infinitesimal automorphism spaces of the coherent stack}
\author{The Stacks Project Authors}
\date{stacks-project:quot.tex:local-snapshot-20261012; GNU FDL 1.2 or later}
\begin{document}
\maketitle
\setcounter{subsection}{1}
\noindent\textit{Editorial note.} All delivered mathematical text is the author's own, copied from the pinned Stacks Project chapter file quot.tex; no statement or proof step was rewritten, reordered or removed. The counted span is the author environment and its own proof environment. References inside the span are handed to the assembler as navigation stubs that carry the canonical tag of the cited tag only; whether each of them is load-bearing is recorded on the stub itself.\par\medskip
\noindent\textit{Allowed standard prerequisites (editorial).}\begin{itemize}
\item Situation situation-coherent (canonical tag 08KB). cited at statement line 1431. Standing situation.
\item Section artin-section-tangent-spaces (canonical tag 07WY). cited at statement line 1438. Definitions of tangent space and Inf.
\item Lemma lemma-coherent-RS-star (canonical tag 08LQ). cited at proof lines 1443-1444. (RS*) of the coherent stack.
\item Section artin-section-inf (canonical tag 07Y6). cited at proof line 1445. Definition of (RS*).
\item Lemma artin-lemma-deformation-category (canonical tag 07WU). cited at proof lines 1448-1449. Fibres are deformation categories.
\item Remark remark-coherent-base-change (canonical tag 08LP). cited at proof line 1456. X\_0's stack as 2-fibre product.
\item Lemma artin-lemma-fibre-product-tangent-spaces (canonical tag 07X2). cited at proof lines 1459-1460. Tangent/Inf of 2-fibre products.
\item Lemma artin-lemma-finite-dimension (canonical tag 07X1). cited at proof lines 1463-1464. Finite dimensionality for B and Spec(k).
\item Lemma lemma-coherent-over-first-order-thickening (canonical tag 08W7). cited at proof lines 1477-1478. Tangent space = lifts of G\_0.
\item Lemma defos-lemma-flat-ringed-topoi (canonical tag 08VW). cited at proof lines 1481 and 1489. Ext\textasciicircum{}1 computes T F; Inf = Ext\textasciicircum{}0.
\item Lemma spaces-perfect-lemma-ext-finite (canonical tag 0D0T). cited at proof lines 1499-1500. Finiteness of the Ext groups.
\end{itemize}\medskip
\ifdefined\hypertarget\hypertarget{sourcepointer-1}{}\fi
\noindent\textit{Source reference.} Original definition 08KB: sources/stacks-chow/tags.txt, source lines 10879--10879. Cited at statement line 1431: the lemma is phrased in this situation; load-bearing.\par\medskip
\ifdefined\hypertarget\hypertarget{sourcepointer-2}{}\fi
\noindent\textit{Source reference.} Original section 07WY: sources/stacks-chow/tags.txt, source lines 10058--10058. Cited at statement line 1438: the objects whose finite dimensionality is claimed; load-bearing.\par\medskip
\ifdefined\hypertarget\hypertarget{sourcepointer-3}{}\fi
\noindent\textit{Source reference.} Original lemma 08LQ: sources/stacks-chow/tags.txt, source lines 10928--10928. Cited at proof lines 1443-1444: X satisfies (RS*), hence (RS); load-bearing.\par\medskip
\ifdefined\hypertarget\hypertarget{sourcepointer-4}{}\fi
\noindent\textit{Source reference.} Original section 07Y6: sources/stacks-chow/tags.txt, source lines 10101--10101. Cited at proof line 1445: definition of (RS*); load-bearing.\par\medskip
\ifdefined\hypertarget\hypertarget{sourcepointer-5}{}\fi
\noindent\textit{Source reference.} Original lemma 07WU: sources/stacks-chow/tags.txt, source lines 10054--10054. Cited at proof lines 1448-1449: the fibres are deformation categories; load-bearing.\par\medskip
\ifdefined\hypertarget\hypertarget{sourcepointer-6}{}\fi
\noindent\textit{Source reference.} Original remark 08LP: sources/stacks-chow/tags.txt, source lines 10927--10927. Cited at proof line 1456: X\_0's stack is the 2-fibre product with Spec(k) over B; load-bearing for the reduction.\par\medskip
\ifdefined\hypertarget\hypertarget{sourcepointer-7}{}\fi
\noindent\textit{Source reference.} Original lemma 07X2: sources/stacks-chow/tags.txt, source lines 10062--10062. Cited at proof lines 1459-1460: reduces to the three factors; load-bearing.\par\medskip
\ifdefined\hypertarget\hypertarget{sourcepointer-8}{}\fi
\noindent\textit{Source reference.} Original lemma 07X1: sources/stacks-chow/tags.txt, source lines 10061--10061. Cited at proof lines 1463-1464: B and Spec(k) have finite dimensional tangent spaces; load-bearing.\par\medskip
\ifdefined\hypertarget\hypertarget{sourcepointer-9}{}\fi
\noindent\textit{Source reference.} Original lemma 08W7: sources/stacks-chow/tags.txt, source lines 11260--11260. Cited at proof lines 1477-1478: the tangent space is the set of lifts of G\_0 to a flat module on X\_epsilon; load-bearing.\par\medskip
\ifdefined\hypertarget\hypertarget{sourcepointer-10}{}\fi
\noindent\textit{Source reference.} Original lemma 08VW: sources/stacks-chow/tags.txt, source lines 11249--11249. Cited at proof lines 1481 and 1489 (twice): T F = Ext\textasciicircum{}1 and Inf = Ext\textasciicircum{}0; load-bearing twice.\par\medskip
\ifdefined\hypertarget\hypertarget{sourcepointer-11}{}\fi
\noindent\textit{Source reference.} Original lemma 0D0T: sources/stacks-chow/tags.txt, source lines 16323--16323. Cited at proof lines 1499-1500: G\_0 coherent on the Noetherian X\_0 gives finite dimensional Exts; load-bearing.\par\medskip
\begin{lemma}
\label{lemma-coherent-tangent-space}
In Situation \ref{situation-coherent} assume that $S$ is a locally Noetherian
scheme and $B \to S$ is locally of finite presentation.
Let $k$ be a finite type field over $S$ and let
$x_0 = (\Spec(k), g_0, \mathcal{G}_0)$
be an object of $\mathcal{X} = \Cohstack_{X/B}$ over $k$. Then
the spaces $T\mathcal{F}_{\mathcal{X}, k, x_0}$ and
$\text{Inf}(\mathcal{F}_{\mathcal{X}, k, x_0})$
(Artin's Axioms, Section \ref{artin-section-tangent-spaces})
are finite dimensional.
\end{lemma}

\begin{proof}
Observe that by Lemma \ref{lemma-coherent-RS-star}
our stack in groupoids $\mathcal{X}$ satisfies property (RS*)
defined in Artin's Axioms, Section \ref{artin-section-inf}.
In particular $\mathcal{X}$ satisfies (RS).
Hence all associated predeformation
categories are deformation categories
(Artin's Axioms, Lemma \ref{artin-lemma-deformation-category})
and the statement makes sense.

\medskip\noindent
In this paragraph we show that we can reduce to the case $B = \Spec(k)$.
Set $X_0 = \Spec(k) \times_{g_0, B} X$
and denote $\mathcal{X}_0 = \Cohstack_{X_0/k}$. In
Remark \ref{remark-coherent-base-change} we have seen that
$\mathcal{X}_0$ is the $2$-fibre product of $\mathcal{X}$ with
$\Spec(k)$ over $B$ as categories fibred in groupoids over
$(\Sch/S)_{fppf}$. Thus by
Artin's Axioms, Lemma \ref{artin-lemma-fibre-product-tangent-spaces}
we reduce to proving that $B$, $\Spec(k)$, and $\mathcal{X}_0$
have finite dimensional tangent spaces and infinitesimal automorphism
spaces. The tangent space of $B$ and $\Spec(k)$ are finite dimensional by
Artin's Axioms, Lemma \ref{artin-lemma-finite-dimension}
and of course these have vanishing $\text{Inf}$.
Thus it suffices to deal with $\mathcal{X}_0$.

\medskip\noindent
Let $k[\epsilon]$ be the dual numbers over $k$.
Let $\Spec(k[\epsilon]) \to B$ be the composition of $g_0 : \Spec(k) \to B$
and the morphism $\Spec(k[\epsilon]) \to \Spec(k)$
coming from the inclusion $k \to k[\epsilon]$.
Set $X_0 = \Spec(k) \times_B X$ and
$X_\epsilon = \Spec(k[\epsilon]) \times_B X$.
Observe that $X_\epsilon$ is a first order thickening of $X_0$
flat over the first order thickening $\Spec(k) \to \Spec(k[\epsilon])$.
Unwinding the definitions and using
Lemma \ref{lemma-coherent-over-first-order-thickening}
we see that $T\mathcal{F}_{\mathcal{X}_0, k, x_0}$ is the set of
lifts of $\mathcal{G}_0$ to a flat module on $X_\epsilon$.
By Deformation Theory, Lemma \ref{defos-lemma-flat-ringed-topoi}
we conclude that
$$
T\mathcal{F}_{\mathcal{X}_0, k, x_0} =
\Ext^1_{\mathcal{O}_{X_0}}(\mathcal{G}_0, \mathcal{G}_0)
$$
Here we have used the identification $\epsilon k[\epsilon] \cong k$
of $k[\epsilon]$-modules. Using
Deformation Theory, Lemma \ref{defos-lemma-flat-ringed-topoi}
once more we see that
$$
\text{Inf}(\mathcal{F}_{\mathcal{X}, k, x_0}) =
\Ext^0_{\mathcal{O}_{X_0}}(\mathcal{G}_0, \mathcal{G}_0)
$$
These spaces are finite dimensional over $k$ as $\mathcal{G}_0$
has support proper over $\Spec(k)$. Namely, $X_0$ is of finite presentation
over $\Spec(k)$, hence Noetherian. Since $\mathcal{G}_0$ is of finite
presentation it is a coherent $\mathcal{O}_{X_0}$-module. Thus we may apply
Derived Categories of Spaces, Lemma
\ref{spaces-perfect-lemma-ext-finite}
to conclude the desired finiteness.
\end{proof}

\end{document}
